\documentclass[10pt,a4paper]{amsart}
\usepackage[margin=19mm]{geometry}
\usepackage{amsmath,amssymb,mathtools}
\usepackage[scr=rsfs,cal=euler]{mathalfa}
\usepackage{xfrac}
\usepackage[unicode,hidelinks,bookmarks=false]{hyperref}
\newcommand{\ie}{i.e.\ }
\newcommand{\cf}{cf.\ }
\newcommand{\resp}{respectively\ }
\newcommand{\Subsection}{\subsection}
\providecommand{\nobreakdash}{\nobreak\hbox{-}}
\setlength{\emergencystretch}{2em}
\multlinegap0pt
\usepackage{mathtools}
\usepackage{xcolor}
\usepackage{float}
\newtheorem{theorem}{Theorem}
\title{Fast-Decaying Polynomial Reproduction}
\author{Stefano De Marchi, Giacomo Cappellazzo}
\date{Source: arXiv:2411.14933v2 (CC BY 4.0)}
\begin{document}
\maketitle

\noindent\emph{Source and licence.} Fast-Decaying Polynomial Reproduction. Version arXiv:2411.14933v2 (CC BY 4.0), distributed by arXiv under the Creative Commons Attribution 4.0 International licence, http://creativecommons.org/licenses/by/4.0/, as recorded in the arXiv licence metadata for this exact version and shown on the abstract page https://arxiv.org/abs/2411.14933v2. Full licence notice: \texttt{licenses/arxiv-2411.14933-CC-BY-4.0.txt}.\par\smallskip

\noindent\textit{Editorial scope.} Counted unit: the article's theorem eq\_stability\_quasi\_interpoaltion\_fast\_decay (source0.tex lines 199--211). The article's complete original proof of it is delivered with it (source0.tex lines 212--263, one \texttt{proof} environment of 52 lines). No mathematics is written, repaired or re-derived: the statement and the proof are the article's own lines, in the article's own notation, and no retained line is substituted or re-flowed.  No further source-local context is needed: every cross-reference of every retained line resolves inside the two delivered spans.  Nothing was omitted from any retained span, so the package carries no omission record.  No retained line contains a citation, so the wrapper carries no bibliography and no bibliography entry.  No retained line contains a figure environment, an image-inclusion command or drawing code, so no image is delivered and the package declares no asset dependency.  Editorial formatting only, in addition: an AMS \texttt{amsart} preamble that restates the article's own declarations and macros used by the retained lines, namely the environments \texttt{theorem} on separate counters, together with the standard packages those lines need.  The retained environments print in this document as Theorem 1 in order of appearance, whereas the article prints the same items with the numbering of its own full document.  The complete native source of the article as distributed in the arXiv e-print for this exact version is bundled unchanged as \texttt{sources/168949/source0.tex}.
\begin{theorem}
    Let $X=\{x_1,\dots,x_N\} \subseteq \mathbb{R}^d$ be an arbitrary data set, and let $\{u_1,\dots,u_N\}$ be given by a fast-decaying polynomial reproduction method with respect to $\varphi$. Then, for every $\ell \in \mathbb{N}$ there exists a constant $K=K(\ell,C,\varphi,d)$ such that 
    \begin{equation}
        \sum_{j=1}^{N} \|x-x_j\|_2^{\ell} |u_j(x)| \leq K q_{X}^{\ell} \quad \textrm{ for } x \in \Omega,
         \label{eq_stability_quasi_interpoaltion_fast_decay}
    \end{equation}
    where
    \begin{equation}
        K=3^d C \sum_{n=0}^{+\infty} (n+1)^{d+\ell-1}\varphi(n).
        \label{def_constant_K_stability}
    \end{equation}
    \label{thm_stability_quasi_interpoaltion_fast_decay}
\end{theorem}

\begin{proof}
Fix $x \in \Omega$. We define a sequence of sets $\{E_n\}_{n \in \mathbb{N}}$ that covers $\mathbb{R}^d$ as
\begin{equation*}
    E_n = \{ y \in \mathbb{R}^d: n q_X \leq \|y-x\|_2 \leq (n+1)q_X\}.
\end{equation*}
We note that if $x_j \in E_n$ then
\begin{equation*}
    B(x_j,q_X) \subseteq \{ y \in \mathbb{R}^d: (n-1)q_X \leq \|y-x\|_2 \leq (n+2)q_X\}=E_{n-1}\cup E_n \cup E_{n+1}\,.
\end{equation*}
In fact, if $y \in B(x_j,q_X)$ then $nq_X -q_X \leq \|x_j-x\|_2-\|x_j-y\|_2 \leq\|y-x\|_2 \leq \|y-x_j\|_2+\|x_j-x\| \leq q_X +(n+1)q_X$.
Since $\{B(x_j,q_X)\}_{j=1,\dots,N}$ are pairwise disjoint,
\begin{equation*}
    \bigcup_{j : x_j \in E_n} B(x_j,q_X) \subseteq \overline{B(x,(n+2)q_X)} \setminus B(x,(n-1)q_X),
\end{equation*}
a volume comparison gives for $n \geq 1$
\begin{equation*}
    \# \{x_j : x_j \in E_n\} \leq (n+2)^d-(n-1)^d \leq 3^d n^{d-1} \leq 3^d(n+1)^{d-1},
\end{equation*}
that holds also for $n=0$ because $2^d \leq 3^d$ ($\#$ being the cardinality of the set).
\\
\\
The foregoing identity is established by induction on the dimension. For $d=1$ we have $n+2-n+1=3$. 
\\ Let us suppose that it is true for $d$, we will prove the inequality for $d+1$.
\begin{equation*}
    \begin{split}
        &(n+2)^{d+1}-(n-1)^{d+1} = (n+2)(n+2)^d-(n-1)(n-1)^d = \\
        = & n((n+2)^d-(n-1)^d)+2(n+2)^d+(n-1)^d = \\
        = & n((n+2)^d-(n-1)^d) +2((n+2)^d-(n-1)^d)+3(n-1)^d \leq\\
        \leq & 3^dn^d+2 \cdot3^dn^{d-1}+3(n-1)^d \leq 3^{d+1}n^d,
    \end{split}
\end{equation*}
where the last inequality holds because by dividing by $3^dn^d$ we have
\begin{equation*}
    1+\frac{2}{n}+ \frac{3}{3^d}\left(\frac{n-1}{n}\right)^d \leq 3,
\end{equation*}
that for $n=1$ becomes $3 \leq 3$, instead for $n\geq 2$ is 
\begin{equation*}
     \frac{2}{n}+ \frac{3}{3^d}\left(\frac{n-1}{n}\right)^d \leq \frac{2}{n}+\left(\frac{n-1}{n}\right)^d \leq \frac{2}{n}+1 \leq 1+1 = 2.
\end{equation*}
Observing that if $x_j \in E_n$ then $n \leq \frac{\|x_j-x\|_2}{q_X}$ and if we split the sum over the sets $\{E_n\}_{n \in \mathbb{N}}$ we obtain 
\begin{equation*}
\begin{split}
    \sum_{j=1}^{N} \|x-x_j\|_2^\ell|u_j(x)| & \leq \sum_{n=0}^{+\infty} \sum_{x_j \in E_n} \|x-x_j\|_2^\ell|u_j(x)| \leq \sum_{n=0}^{+\infty}\sum_{x_j \in E_n} \|x-x_j\|_2^\ell C\varphi\left( \frac{\|x-x_j\|_2}{q_X} \right) \leq \\
    & \leq \sum_{n=0}^{+\infty}\sum_{x_j \in E_n} \|x-x_j\|_2^\ell C\varphi(n) \leq \sum_{n=0}^{+\infty} \sum_{x_j \in E_n} (n+1)^\ell q_X^\ell C\varphi(n) \leq \\
    & \leq \sum_{n=0}^{+\infty} 3^d (n+1)^{d+\ell-1}q_X^\ell C \varphi(n) = q_{X}^{\ell} \underbrace{3^d C \sum_{n=0}^{+\infty} (n+1)^{d+\ell-1}\varphi(n)}_{K}.
\end{split}
\end{equation*}
Convergence of the series follows from the ratio test:
\begin{equation*}
    \frac{(n+2)^{d+\ell-1} \varphi(n+1)}{(n+1)^{d+\ell-1} \varphi(n)} = \left( \frac{n+2}{n+1}\right)^{d+\ell-1} \frac{\varphi(n+1)}{\varphi(n)} \xrightarrow[]{n \rightarrow +\infty}  \lim_{n\to +\infty} \frac{\varphi(n+1)}{\varphi(n)} < 1.
\end{equation*}
\end{proof}

\end{document}
